% !TEX program = lualatex
\documentclass[a4paper,11pt]{ltjsarticle}

\usepackage{amsmath,amssymb}
\usepackage{booktabs}
\usepackage{geometry}
\usepackage{graphicx}
\usepackage{tikz}
\usepackage{xcolor}
\usepackage{enumitem}
\usepackage{hyperref}
\geometry{margin=24mm}
\usetikzlibrary{arrows.meta,positioning,shapes.geometric,calc}

\title{進捗報告メモ\\多目的GPにおけるBO制御実装の改善}
\author{}
\date{2026年5月26日}

\newcommand{\HV}{\mathrm{HV}}
\newcommand{\ND}{\mathrm{ND}}

\begin{document}
\maketitle

\section{研究の狙い}

本研究では，多目的遺伝的プログラミング (MOGP) において，
交叉率 \(p_c\)，突然変異率 \(p_m\)，更新周期 \(k\) を固定せず，
現在の探索状態に応じてベイズ最適化 (BO) で逐次決定する手法を検討している．
目的は，最終的に良い Pareto front を得ることと，探索中の多様性を維持することである．

\begin{figure}[htbp]
\centering
\resizebox{0.93\linewidth}{!}{%
\begin{tikzpicture}[
  node distance=10mm and 14mm,
  block/.style={rectangle, rounded corners=2pt, draw, align=center, minimum width=38mm, minimum height=9mm, fill=blue!6},
  bo/.style={rectangle, rounded corners=2pt, draw, align=center, minimum width=38mm, minimum height=9mm, fill=orange!12},
  gp/.style={rectangle, rounded corners=2pt, draw, align=center, minimum width=38mm, minimum height=9mm, fill=green!10},
  arrow/.style={-Latex, thick}
]
\node[block] (state) {GP状態\\ HV, 多様性, 木サイズ};
\node[bo, right=of state] (controller) {BO制御器\\ \(p_c,p_m,k\) を提案};
\node[gp, below=of controller] (plant) {多目的GP\\ \(k\) 世代実行};
\node[block, left=of plant] (reward) {報酬計算\\ HV改善 + 多様性};
\draw[arrow] (state) -- (controller);
\draw[arrow] (controller) -- (plant);
\draw[arrow] (plant) -- (reward);
\draw[arrow] (reward) -- (state);
\end{tikzpicture}
}
\caption{提案手法の基本構造}
\end{figure}

\section{実装改善の流れ}

ここまでの実装は，次の 3 段階で改善してきた．

\begin{table}[htbp]
\centering
\caption{実装バージョンの整理}
\begin{tabular}{lll}
\toprule
版 & 目的 & 改善内容 \\
\midrule
ver1 & まず動かす & BO 制御 GP と普通 GP を比較 \\
ver2 & 評価を正確にする & archive の重複除去と世代別ログを追加 \\
ver2\_b & 本格実験に備える & archive 評価と population 観測を分離 \\
\bottomrule
\end{tabular}
\end{table}

ver1 では，提案手法が実際に動き，普通 GP より良い傾向があるかを確認した．
しかし，archive に重複が含まれ，候補解数を過大に見積もる可能性があった．

そこで ver2 では，同じ目的値・同じ構造の個体を重複除去し，
unique objective，unique structure，unique pair を記録するようにした．

さらに ver2\_b では，次のように役割を分けた．
\[
\begin{array}{ll}
\text{archive} &: \text{HV，報酬，最終 Pareto front の評価}\\
\text{population} &: \text{現在の多様性，平均木サイズなどの状態観測}
\end{array}
\]

archive 更新は
\[
A_{g+1}=\ND(A_g\cup P_g\cup Q_g)
\]
として，offspring に一度現れた良い非劣解も保存できるようにした．

\section{主な結果}

\subsection{テンプレート問題}

テンプレート問題では，提案手法が普通 GP より高い HV と Diversity を示した．
ver2\_b では，提案手法の最終 archive HV は 0.989710，
普通 GP は 0.979405 であった．
また，Diversity も提案手法の方が高かった．

\begin{table}[htbp]
\centering
\caption{テンプレート問題の最終結果}
\begin{tabular}{llll}
\toprule
手法 & archive HV & Diversity & archive size \\
\midrule
提案手法 ver2\_b & 0.989710 & 0.592879 & 6 \\
普通 GP & 0.979405 & 0.347826 & 2 \\
\bottomrule
\end{tabular}
\end{table}

\begin{figure}[htbp]
\centering
\includegraphics[width=0.88\linewidth]{outputs/template_plain_vs_bogp_test_ver2_b1/20260520_173956/archive_population_hv_diversity_comparison.png}
\caption{テンプレート問題: archive HV，population HV，population diversity}
\end{figure}

\subsection{構造探索問題}

構造探索問題でも，提案手法は普通 GP よりわずかに高い HV と Diversity を示した．
ただし，HV は早期に高い値へ到達しており，問題設定としては少し簡単すぎる可能性がある．

\begin{table}[htbp]
\centering
\caption{構造探索問題の最終結果}
\begin{tabular}{lllll}
\toprule
手法 & archive HV & Diversity & archive size & unique structure \\
\midrule
提案手法 ver2\_b2 & 0.999546 & 0.783334 & 10 & 10 \\
普通 GP & 0.999533 & 0.750047 & 7 & 7 \\
\bottomrule
\end{tabular}
\end{table}

\begin{figure}[htbp]
\centering
\includegraphics[width=0.88\linewidth]{outputs/structural_plain_vs_bogp_test_ver2_b2/20260520_172026/archive_population_hv_diversity_comparison.png}
\caption{構造探索問題: archive HV，population HV，population diversity}
\end{figure}

\section{分かったこと}

\begin{enumerate}[leftmargin=3em]
  \item 提案手法は，小規模実験では普通 GP より HV と Diversity を高く保つ傾向がある．
  \item ver1 の archive size は重複を含んでおり，実質的な Pareto 候補数としては過大であった．
  \item ver2 以降では，unique objective / structure / pair を記録することで，結果をより正確に説明できるようになった．
  \item ver2\_b では，archive を成果評価，population を状態観測に使う整理ができた．
  \item 構造探索では，パラメータ値込みの構造キーを使うと，同じ目的値でも異なる設計候補を保持できる．
\end{enumerate}

\section{今後の課題}

今後は，小規模な動作確認から本格実験へ進む．
特に次の点が重要である．

\begin{itemize}[leftmargin=3em]
  \item 複数 seed で実験し，平均・標準偏差を確認する．
  \item 構造探索問題の HV 参照点や正規化範囲を見直す．
  \item 比較対象を固定率 GP だけでなく，複数固定率，手設計スケジュール，非文脈 BO に広げる．
  \item 多様性項なし，\(k\) 制御なし，構造キー b1/b2 などの ablation を行う．
  \item 最終的には HV，Diversity，unique objective，unique structure，制御回数をまとめて評価する．
\end{itemize}

\section{まとめ}

ここまでで，提案手法の最小実装から，評価基盤を整えた ver2\_b まで進めた．
現時点では，提案手法は普通 GP より良い傾向を示しているが，
構造探索問題では HV が飽和しやすく，より厳密な比較には本格実験が必要である．
次の段階では，複数 seed と比較手法を揃え，提案手法の有効性を検証していく．

\end{document}
